\section{Discussion}
\subsection{Attribution, Adaptation, and Operational Scale}
The primary learned energy reduction is mainly detector choice: both learned policies execute LightLR, and the matched control resolves no Tabular-Q saving or premium against direct LightLR. Its 13.1 J contrast with Static-MedRF is only 26.1 mW on average while missed-positive presentations increase almost tenfold. At the primary 100 rows/s cadence, detector inference occupies only a small fraction of each one-second interval, so background platform power dominates the whole-device endpoint. Energy alone would therefore overstate the practical value of the learned policy.

The variable-load campaign supplies the positive service-management result that the primary campaign lacked. Relative to Static-MedRF, DeadlineGuard reduces energy by 43.7 J per run (2.98\%; 95\% paired CI 41.4--46.0 J) and eliminates a 10.24-pp deadline-miss rate, at a 1.373-pp F1 cost. Unlike the learned selectors, it realizes load-conditioned switching under the prospectively frozen factorial. This does not make DeadlineGuard universally optimal: Static-LightLR is still lower-energy and has the same zero deadline misses, while DeadlineGuard trades +7.55 J for a small 0.056-pp F1 gain relative to LightLR. The value of the comparison is that the high-security static baseline becomes infeasible at the declared service deadline and a simple causal rule exposes that trade-off explicitly.

The confidence cascade answers a different question. It recovers MedRF-like F1 using only 1.34\% row escalation, but those rows occur in 66.76\% of 100-row windows. The executed implementation performs one MedRF call per escalation-active window on a tiny deferred sub-batch, and that call has near-MedRF latency. Its observed whole-device cost is therefore specific to this call granularity and should not be read as a general limit on optimized batched cascades.

\subsection{Reward Specification as a Pre-deployment Check}
The policy collapse is not only an empirical observation. Equation~\eqref{eq:nondeg} exposes whether the frozen scalar utility even permits a pair of actions to exchange rank on the declared support. For MedRF versus LightLR, candidate-range normalization implies a 0.7967 balanced-accuracy threshold---an almost 80-percentage-point exchange rate determined before held-out evaluation. A controller can therefore appear to ``learn'' a stable choice even when the declared objective makes reversal implausibly demanding on its calibration/validation support. The broader design lesson is to normalize resource terms against deployment contracts---idle-relative power, latency SLOs, energy budgets, or thermal headroom---rather than the arbitrary extrema of the current candidate pool.

\subsection{Sparse Support and Off-support Safety}
Tabular-Q's full greedy map assigns 314/324 states to TinyDT and only ten to LightLR, while formal primary execution reaches four LightLR states. Frozen source confirms that this is not a separate runtime-exception fallback: unseen zero-initialized Q rows are tied, \texttt{np.argmax} returns the first maximum, and TinyDT is action index zero. variable-load campaign makes the consequence concrete: new load-related states drive the frozen table to TinyDT, increasing FPR by roughly 66 percentage points versus LightLR. The transferable risk is therefore the interaction of sparse support, unsafe default action ordering, and deterministic tie breaking. DQN maps the primary 324-state grid to LightLR, but that smoother off-support behavior is not evidence of validated generalization. A deployable discrete controller should report state coverage, the full action map, and the safety quality of unvisited/default actions---not only occupancy on visited states.

\subsection{Measurement and Evaluation Implications}
The variable-load campaign also illustrates why raw means alone are insufficient at small whole-device effects. DQN executes LightLR in every window, yet its raw energy difference from Static-LightLR is positive in 10/10 blocks while the Student-$t$ interval crosses zero because one thermal/power excursion dominates the variance. Reporting the block points, SD, median, exact sign-flip sensitivity, and preceding-idle adjustment exposes this instability; none licenses a causal controller-power explanation because the controller-only power study stopped before measurement.

A learned orchestrator should therefore be compared with the static action it realizes and with simple conditional controls when workload feasibility changes. Action occupancy, deadline violations, and confusion counts are indispensable: energy alone would favor the failed TinyDT collapse, while F1 alone would hide false alarms and service infeasibility. Raw and idle-adjusted physical endpoints answer different questions and neither is selected post hoc because it looks favorable.

The TinyDT diagnostic shows that its threshold-0.5 failure is substantially calibration-sensitive. That does not legitimize replacing the primary threshold after test inspection or treating TinyDT as safe in the already frozen physical campaigns. It instead motivates validation-only calibration governance before deployment. Generalization beyond this operating envelope still requires untouched traffic, additional boards, and independent metrology.
